\documentclass[12pt,a4paper]{article}
\usepackage[utf8]{inputenc}
\usepackage[T1]{fontenc}
\usepackage[polish]{babel}
\usepackage{amsmath,booktabs,longtable,tabularx}
\usepackage{geometry,enumitem,xcolor,hyperref,fancyhdr}
\geometry{margin=2.5cm}
\hypersetup{colorlinks=true,linkcolor=black,citecolor=blue!60!black,
urlcolor=blue!60!black,pdftitle={Ocena kompletności i wiarygodności danych},
pdfauthor={Lidia Cichońska}}
\setlist{itemsep=0.25em,topsep=0.4em}
\renewcommand{\arraystretch}{1.2}
\sloppy
\pagestyle{fancy}
\fancyhf{}
\lhead{Wiarygodność i kompletność danych w HCI}
\rhead{Lidia Cichońska}
\cfoot{\thepage}
\newcommand{\pp}{p.p.}

\title{Propozycja oceny kompletności, reprezentatywności\\
i wiarygodności danych w systemach analizy obrazu}
\author{Lidia Cichońska}
\date{25 sierpnia 2026 r.}

\begin{document}
\maketitle

\begin{abstract}
System HCI jest wiarygodny tylko w takim zakresie, w jakim dane użyte do jego
zaprojektowania i oceny opisują zjawisko, o którym system ma wnioskować. W
analizie obrazu są to nie tylko obrazy i etykiety, lecz również podziały
zbioru, rejestry predykcji, wartości pewności, reprezentacje ukryte, decyzje o
wykonaniu części sieci oraz pomiary czasu i kosztu.

W sprawozdaniu zaproponowano \emph{Mapę Pokrycia i Wiarygodności Danych}
(MPWD) dla specjalisty projektującego lub audytującego system. Ma ona
pokazywać, które fragmenty zjawiska są dobrze opisane, gdzie brakuje danych i
w jakim zakresie wolno uogólniać wynik. Przeprowadzone eksperymenty z
klasyfikacją obrazów i adaptacyjnym wykonywaniem sieci są studium przypadku,
a nie pełnym raportem z badań.
\end{abstract}

\section{Punkt wyjścia i rozwinięcie wcześniejszej propozycji}

Niniejsze sprawozdanie jest bezpośrednim rozwinięciem wcześniejszej propozycji
oceny wiarygodności channel gatingu. W tamtym ujęciu punktem wyjścia było
pytanie techniczne: czy deklarowana kompresja modelu jest wiarygodna, jeżeli
średnia accuracy pozostaje wysoka. Odpowiedzią była \emph{decision map}
zbudowana z markerów M1--M6: jakości per klasa, koncentracji błędów,
stabilności progu, zależności budżetu od klasy, zmiany reprezentacji i
odporności na kwantyzację.

Podejście to było użyteczne, ale zbyt szybko przechodziło od danych do
szczegółów mechanizmu kompresji. Nie wyjaśniało dostatecznie wyraźnie, jakie
zjawisko ma być reprezentowane, kto korzysta z oceny i kiedy wynik wolno
uogólnić poza benchmark. Obecna propozycja nie odrzuca wcześniejszych
markerów. Umieszcza je wewnątrz szerszej oceny kompletności danych i zakresu
stosowalności systemu.

\begin{quote}
Wcześniejsza praca pytała, czy wynik kompresji modelu jest wiarygodny.
Obecna praca pyta wcześniej i szerzej: czy dane wystarczająco opisują
zjawisko, aby wynik działania i kompresji modelu mógł być wiarygodny dla
określonego zastosowania HCI.
\end{quote}

\begin{table}[ht]
\centering
\caption{Ciągłość pomiędzy wcześniejszą decision map a obecną propozycją.}
\begin{tabularx}{\textwidth}{p{4.5cm}p{4.0cm}X}
\toprule
Wcześniejszy marker i jego nazwa & Rola w obecnej MPWD & Rozwinięcie wynikające z badań \\
\midrule
M1 --- zachowanie jakości w klasach; M2 --- koncentracja błędów & D5: równość jakości i kosztu & Analiza nie kończy się na klasach;
obejmuje również regiony przestrzeni ukrytej i przypadki słabo pokryte. \\
M3 --- stabilność decyzji progowej & D4: stabilność wniosków & Wrażliwość progu jest oceniana razem z
niepewnością, seedami i przesunięciem rozkładu. \\
M4 --- zależność budżetu od klasy & D5: równość jakości i kosztu & Budżet jest analizowany per klasa lub
region danych, a nie wyłącznie jako średnia całego zbioru. \\
M5 --- zmiana reprezentacji wewnętrznej & D3: pokrycie przestrzeni reprezentacji & Reprezentacja nie służy już
tylko do porównania modeli; pomaga wykrywać luki w danych. \\
M6 --- odporność aktywacji na kwantyzację & D6: adekwatność do zastosowania & Odporność techniczna jest jednym z
warunków wdrożenia, obok zgodności domeny, urządzenia i warunków użycia. \\
Decision map --- mapa decyzji & Raport MPWD dla specjalisty & Zachowano statusy
``wystarczający'', ``ostrzeżenie'' i ``brak danych'', ale dodano granicę
stosowalności oraz informację, jakich danych brakuje. \\
\bottomrule
\end{tabularx}
\end{table}
W dalszej części symbole D1--D6 zawsze oznaczają nowe markery danych, natomiast M1--M6 odnoszą się wyłącznie do markerów z wcześniejszego sprawozdania.

Kolejne eksperymenty dostarczyły podstaw do rozszerzenia. Audyt duplikatów
pokazał, że wiarygodność zaczyna się przed oceną modelu. Wrażliwość bramek na
próg wykazała, że pojedynczy parametr nie wystarcza do stabilnego wniosku.
Eksperymenty wieloseedowe pokazały potrzebę raportowania niepewności, a
rzeczywisty pomiar prefix-forward ujawnił różnicę pomiędzy wskaźnikiem
pośrednim i korzyścią możliwą do odczucia w systemie. MPWD scala te obserwacje
w jeden mechanizm oceny danych.

\section{Dziedzina, zastosowanie i użytkownik}

HCI występuje wszędzie tam, gdzie człowiek korzysta z systemu informatycznego
lub opiera na nim decyzję. W rozpatrywanym obszarze system analizuje obraz,
wyznacza predykcję i może dostosować ilość obliczeń do wejścia. Może być
częścią aplikacji mobilnej, narzędzia wspomagającego operatora albo komponentu
na urządzeniu o ograniczonych zasobach.

Użytkownikiem proponowanej oceny jest specjalista: badacz, inżynier uczenia
maszynowego, właściciel danych lub audytor. Powinien otrzymać odpowiedź:

\begin{enumerate}
 \item czy dane opisują zjawisko, które system ma analizować;
 \item gdzie wyniki systemu są wiarygodne;
 \item gdzie dane są niepełne, stronnicze lub odmienne od przyszłego użycia;
 \item o jakich ograniczeniach trzeba poinformować odbiorcę systemu.
\end{enumerate}

System może być wiarygodny w wąskiej poddziedzinie, nie będąc uniwersalnym.
Wynik na małych obrazach benchmarkowych może rzetelnie porównywać metody w tym
benchmarku, ale nie uzasadnia użycia modelu dla obrazów medycznych, zdjęć z
telefonu ani innej domeny. Granica stosowalności nie jest wadą, jeżeli zostaje
jawnie zakomunikowana.

\section{Zjawisko i opisujące je dane}

Badane zjawisko to zależność między treścią obrazu, trudnością klasyfikacji,
decyzją modelu i kosztem jej uzyskania. Interesuje nas nie tylko średnia
trafność, lecz także stabilność jakości i kosztu w klasach i różnych obszarach
przestrzeni obrazów.

W eksperymentach wykorzystano publiczne benchmarki CIFAR-10, CIFAR-100 i Tiny
ImageNet. Oprócz obrazów i etykiet analiza wykorzystuje dane systemowe:
identyfikator i podział próbki, predykcję, pewność, informację o błędzie,
reprezentację ukrytą, wybrane wyjście sieci, koszt, czas wykonania, seed i
konfigurację. Połączenie danych z obrazu, bazy i rejestrów systemowych pozwala
sprawdzić, czy oszczędność obliczeń nie ukrywa pogorszenia jakości.

Przyjęto następujące znaczenia pojęć:

\begin{description}
 \item[Wiarygodność] znane pochodzenie danych, kontrola etykiet i podziałów
 oraz możliwość odtworzenia wyniku.
 \item[Kompletność] obecność obserwacji potrzebnych do opisu istotnych
 wariantów zjawiska i wyznaczenia deklarowanych miar.
 \item[Reprezentatywność] zgodność danych badawczych z populacją lub
 kontekstem, do którego odnosi się wniosek.
 \item[Rzetelność oceny] stabilność wniosku względem próby, seeda, rozsądnej
 zmiany progu i wyboru miary, wraz z pokazaniem niepewności.
\end{description}

Nie istnieją dane reprezentatywne ``w ogóle''. Ten sam zbiór może być dobry
do porównania metod i niewystarczający do oceny systemu w realnym środowisku.

\section{Mapa Pokrycia i Wiarygodności Danych}

MPWD odpowiada na pytanie: \emph{czy mamy dobre dane do tego, co chcemy
analizować?} Powstaje, gdy dostępne są obrazy, predykcje, reprezentacje ukryte
i rejestry kosztu. Kryteria powinny być jednak ustalone przed końcową analizą,
aby ograniczyć dobieranie miar do korzystnego wyniku. MPWD jest kartą oceny,
a nie jedną liczbą.

\begin{longtable}{p{1.0cm}p{3.6cm}p{4.4cm}p{5.2cm}}
\caption{Markery Mapy Pokrycia i Wiarygodności Danych.}\\
\toprule
Kod & Nazwa & Pytanie & Sposób oceny \\
\midrule
\endfirsthead
\toprule
Kod & Nazwa & Pytanie & Sposób oceny \\
\midrule
\endhead
D1 & Integralność i niezależność danych & Czy próbki, etykiety i podziały są
poprawne? & Braki, duplikaty, konflikty etykiet, przecieki train--test i
pochodzenie danych. \\
D2 & Pokrycie zjawiska & Czy obecne są istotne warianty zjawiska? &
Liczebności klas i kontekstów, kombinacje cech oraz mało liczne i puste
obszary. \\
D3 & Pokrycie przestrzeni reprezentacji & Czy dane obejmują różnorodność
widoczną dla modelu? & Mapa reprezentacji ukrytej, lokalna gęstość, odległość
od danych uczących i obserwacje odstające. \\
D4 & Stabilność wniosków & Czy wynik utrzymuje się po zmianie próby lub
ustawień? & Bootstrapowe przedziały ufności, wiele seedów, wrażliwość progów i
przesunięcie rozkładu. \\
D5 & Równość jakości i kosztu & Czy część danych nie jest systematycznie gorzej
obsługiwana? & Trafność, kalibracja, dodatkowe błędy i koszt per klasa lub per
region danych. \\
D6 & Adekwatność do zastosowania & Czy dane odpowiadają warunkom planowanego
użycia? & Porównanie źródła, rozdzielczości, klas, urządzenia i kontekstu z
docelowym zastosowaniem; lista wyłączeń. \\
\bottomrule
\end{longtable}

Markery dotyczą specyfiki zjawiska, nie tylko ogólnego próbkowania. D2 ocenia
odmiany obrazu istotne dla rozpoznawania, D3 ujawnia luki widoczne dla modelu,
a D5 łączy dane wejściowe z jakością i kosztem działania systemu.

\subsection{Przestrzeń ukryta jako mapa braków}

Etykiety klas nie opisują całej zmienności. Obrazy tej samej klasy mogą różnić
się tłem, jakością, perspektywą lub trudnością. Dla każdego obrazu proponuje
się zapisać reprezentację $h(x)$ i wyznaczyć średnią odległość od $k$
najbliższych danych uczących:

\begin{equation}
 d_k(x)=\frac{1}{k}\sum_{j\in N_k(x)}\lVert h(x)-h(x_j)\rVert_2.
\end{equation}

Duża wartość $d_k(x)$ wskazuje obszar słabo pokryty. Dla specjalisty można
przygotować mapę PCA lub UMAP, kolorując punkty według klasy, błędu, pewności i
kosztu. Wizualizacja wskazuje obszary do kontroli; właściwe miary należy liczyć
również w oryginalnej przestrzeni. Ponieważ reprezentację tworzy model i może
ona powielać jego obciążenia, mapę trzeba konfrontować z opisem wejść,
etykietami i prostym modelem odniesienia.

\subsection{Wynik dla specjalisty}

Końcowy raport MPWD powinien zawierać cel i badaną populację, status każdego
markera (wystarczający, ostrzegawczy lub niewystarczający), mapę dobrze i słabo
pokrytych regionów, jakość i koszt w tych regionach, przedziały ufności oraz
jawną granicę stosowalności. Nie należy uśredniać markerów. Przeciek danych w
D1 może blokować wniosek mimo dobrych pozostałych miar. Ograniczony D6 nie
unieważnia porównania benchmarkowego, lecz blokuje uogólnienie na inną domenę.

\section{Studium przypadku z przeprowadzonych badań}

Poniższe obserwacje są przykładami użycia markerów, a nie kroniką badań.

\begin{table}[ht]
\centering
\caption{Przykłady zastosowania markerów.}
\begin{tabularx}{\textwidth}{p{1.0cm}p{4.2cm}X}
\toprule
Marker & Obserwacja & Interpretacja \\
\midrule
D1 & W CIFAR-10 nie wykryto dokładnego nakładania train--test. W lokalnej
kopii CIFAR-100 znaleziono 10 identycznych obrazów po obu stronach podziału. &
Wniosek dla CIFAR-100 wymaga raportowania i analizy po usunięciu próbek. \\
D4 & Udział małych bramek silnie zależał od progu; fragile ratio osiągał
99,9--100\%. & Małe wartości bramek nie są stabilnym dowodem rzeczywistej
redukcji kosztu. \\
D4/D5 & W trzech seedach adaptacyjne wyjście zmieniło średnią trafność o
$-0{,}07\pm0{,}21$ \pp{} względem pełnej sieci przy koszcie 84,92\%. & Wynik
jest powtarzalnym sygnałem, ale wymaga oceny per klasa i region reprezentacji. \\
D5 & Prefix-forward dla batch size 1 wskazał około 14,6\% krótszą oczekiwaną
medianę czasu. & Pominięcie obliczeń potwierdzono tylko dla zbadanej
konfiguracji, nie dla innych urządzeń i sposobów batchowania. \\
D6 & CIFAR i Tiny ImageNet są małymi benchmarkami. & Umożliwiają porównanie
metod, ale nie reprezentują dowolnych obrazów ani populacji użytkowników. \\
\bottomrule
\end{tabularx}
\end{table}

Trzeba rozdzielić trzy twierdzenia: model generuje małe wartości bramek, graf
faktycznie pomija operatory, a pominięcie daje korzyść na docelowym urządzeniu.
Każde wymaga innych danych. Parametry nie zastępują rejestru operatorów, a
liczba operatorów nie zastępuje pomiaru opóźnienia w kontekście użycia.

Obecne dane wystarczają do kontrolowanego porównania wariantów i stwierdzenia,
że wczesne wyjście pomija część obliczeń dla batch size 1 w zbadanym
środowisku. Nie wystarczają do deklaracji tej samej jakości i korzyści dla
dowolnego obrazu, urządzenia i kontekstu HCI. Komunikat dla specjalisty:

\begin{quote}
System jest wiarygodnie oceniony jako prototyp benchmarkowy dla małych zbiorów
klasyfikacji obrazów. Mocną stroną jest zachowanie średniej jakości i
potwierdzone pomijanie części obliczeń. Luka dotyczy pokrycia przestrzeni
danych, wyników w klasach, przesunięcia domeny i innych urządzeń. Nie należy
jeszcze uogólniać wyniku na zastosowania o innym rozkładzie obrazów.
\end{quote}

\section{Proponowana analiza uzupełniająca}

Pierwszą MPWD można przygotować bez nowego treningu. Dla każdego obrazu należy
zapisać etykietę, predykcję pełną i adaptacyjną, pewność, wybrane wyjście,
koszt oraz reprezentację ukrytą. Następnie należy:

\begin{enumerate}
 \item powtórzyć audyt duplikatów i analizę po ich wykluczeniu;
 \item policzyć liczebność, trafność, kalibrację i koszt per klasa;
 \item zbudować mapę przestrzeni ukrytej i wskazać regiony małej gęstości;
 \item porównać błędy i koszt w regionach dobrze i słabo pokrytych;
 \item wyznaczyć bootstrapowe 95\% przedziały ufności;
 \item zbadać stabilność po niewielkim przesunięciu progów;
 \item wygenerować kartę zakresu stosowalności według D1--D6.
\end{enumerate}

Zbyt mała liczebność regionu powinna dawać status ``nierozstrzygnięty'', a nie
być traktowana jako dowód równości. W kolejnej iteracji można dodać szum,
rozmycie i zmianę jakości obrazu, aby sprawdzić, czy marker przestrzeni ukrytej
przewiduje spadek jakości.

\section{Miejsce oceny i odpowiedzialność}

Ocena może wystąpić przed treningiem (pochodzenie, braki, duplikaty), po
treningu (przestrzeń ukryta, błędy i kalibracja), po wyborze polityki
adaptacyjnej (stabilność progu i koszt), przed wdrożeniem (zgodność z
docelowym środowiskiem i latency) oraz podczas eksploatacji (odległość nowych
danych od znanej przestrzeni).

Ostatnim momentem przed przekazaniem wyniku człowiekowi jest generowanie
odpowiedzi. Dla wejścia odległego od znanych danych system mógłby dołączyć
ostrzeżenie, skierować przypadek do specjalisty albo uruchomić pełną ścieżkę.
W ramach zadania narzędzie nie będzie implementowane; propozycja określa dane
i kryteria potrzebne do jego późniejszego zbudowania.

Klasy benchmarkowe nie są grupami społecznymi, dlatego analiza per klasa nie
jest pełnym badaniem fairness. Jest próbą metodologiczną wykrywania
nierównomiernego pokrycia. W zastosowaniu dotyczącym ludzi markery trzeba
zdefiniować z ekspertem dziedzinowym, uwzględniając cechy chronione,
konsekwencje błędu i możliwość zakwestionowania decyzji. MPWD nie może tworzyć
pozoru obiektywności przez sprowadzenie problemów do jednej liczby.

\section{Wniosek}

Wiarygodność systemu analizy obrazu zależy od zgodności danych ze zjawiskiem,
integralności podziałów, pokrycia wariantów obrazu, stabilności wyniku oraz od
tego, czy jakość i koszt nie pogarszają się w słabo reprezentowanych częściach
danych.

MPWD porządkuje ocenę za pomocą sześciu markerów: integralności i
niezależności danych, pokrycia zjawiska, pokrycia przestrzeni reprezentacji,
stabilności wniosków, równości jakości i kosztu oraz adekwatności do
zastosowania. Jej wynikiem nie jest etykieta ``model wiarygodny'', lecz
informacja, gdzie system jest mocny, gdzie brakuje danych i poza jakim
zakresem należy ostrzec przed uogólnieniem. Dostępne eksperymenty pozwalają
przygotować pierwszą taką ocenę bez ponownego treningu.

\begin{thebibliography}{9}
\bibitem{cifar} A. Krizhevsky, \emph{Learning Multiple Layers of Features
from Tiny Images}, 2009.
\bibitem{resnet} K. He i in., ``Deep Residual Learning for Image
Recognition'', CVPR, 2016.
\bibitem{senet} J. Hu, L. Shen, G. Sun, ``Squeeze-and-Excitation Networks'',
CVPR, 2018.
\bibitem{branchynet} S. Teerapittayanon, B. McDanel, H. T. Kung,
``BranchyNet: Fast Inference via Early Exiting from Deep Neural Networks'',
ICPR, 2016.
\bibitem{datasetshift} J. Qui\~nonero-Candela i in. (red.), \emph{Dataset
Shift in Machine Learning}, MIT Press, 2009.
\bibitem{modelcards} M. Mitchell i in., ``Model Cards for Model Reporting'',
FAT*, 2019.
\end{thebibliography}

\end{document}
